% Irodov Part III §3.7 original 372--388, source PDF pp.156--158
% (printed 147--149). Detailed derivations and checks; original numbering.
\electrochargedproblem{372}{147}
平行板电容器极板距离 $\ell=5.0\,\mathrm{cm}$。一电子在 $t=0$ 从一极板由静止出发，极板间加速电压按 $V(t)=at$ 增加，其中 $a=100\,\mathrm{V/s}$。求电子抵达另一极板时的速度。
\exsolution
板间场强 $E=at/\ell$，电子加速度模
$\ddot x=eat/(m_e\ell)$。以 $x(0)=\dot x(0)=0$ 两次积分：
\[
v(t)=\frac{ea}{2m_e\ell}t^2,\qquad
x(t)=\frac{ea}{6m_e\ell}t^3.
\]
到达另一极板 $x(t_*)=\ell$，有
$t_*=[6m_e\ell^2/(ea)]^{1/3}$。因此
\[
\boxed{v_*=\frac{ea}{2m_e\ell}
\left(\frac{6m_e\ell^2}{ea}\right)^{2/3}
\simeq1.6\times10^4\,\mathrm{m/s}=16\,\mathrm{km/s}}.
\]
\excheck 电子速度远低于光速，非相对论模型成立。

\electrochargedproblem{373}{147}
质子先经电压 $V$ 加速，沿平行板电容器的板长方向射入，板长为 $\ell$，横向匀强电场按 $E(t)=at$ 增强，质子在 $t=0$ 进入场区。忽略边缘效应，求进出场时速度方向夹角 $\theta$。
\exsolution
非相对论质子初始前向速度
$v_x=\sqrt{2eV/m_p}$，
穿越时间 $t_*=\ell/v_x$。
横向加速度
$\dot v_y=ea t/m_p$，
积分得
$v_y(t_*)=ea t_*^2/(2m_p)
=ea\ell^2/(2m_pv_x^2)$。
因此
\[
\boxed{\tan\theta=\frac{v_y}{v_x}
=\frac{ea\ell^2}{2m_p(2eV/m_p)^{3/2}}
=\frac{a\ell^2}{4V}\sqrt{\frac{m_p}{2eV}}}.
\]
\excheck 横向电场不对纵向速度做功，$v_x$ 恒定；$\tan\theta$ 维度必须为无量纲。

\electrochargedproblem{374}{147}
比荷为 $q/m$ 的粒子在一维电场
$E(x)=E_0-ax$ 中由原点静止出发，
$a>0$，假设初始力指向 $+x$。
求 (a) 粒子再次停下前位移；(b) 那一刻加速度。
\exsolution
功—能定理：
$\frac12mv(x)^2=q\int_0^x(E_0-as)\,\mathrm ds
=q(E_0x-ax^2/2)$。
除初始位置 $x=0$ 外，
再次停下满足
$E_0x-ax^2/2=0$，
所以
\[
\boxed{x_*=\frac{2E_0}{a}}.
\]
此时场强
$E(x_*)=E_0-2E_0=-E_0$，
故有向加速度
\[
\boxed{\ddot x_*=-\frac{qE_0}{m}},
\quad |\ddot x_*|=\frac{|qE_0|}{m}.
\]
\excheck 粒子到达停止点前经过 $x=E_0/a$ 时电场改变方向，但速度直到 $x_*$ 才再次降为零。

\electrochargedproblem{375}{147}
电子从静止开始，在均匀电场
$E=10\,\mathrm{kV/cm}$ 中运动。
求动能达到其静止能量 $m_ec^2$ 的时间。
\exsolution
恒定电场作用下动量满足
$p=eEt$。
相对论动能
$T=\sqrt{m_e^2c^4+p^2c^2}-m_ec^2$。
令 $T=m_ec^2$，得
$\sqrt{m_e^2c^4+p^2c^2}=2m_ec^2$，
故 $p=\sqrt3m_ec$。
于是
\[
\boxed{t=\frac{\sqrt3m_ec}{eE}}.
\]
$E=10\,\mathrm{kV/cm}=1.0\times10^6\,\mathrm{V/m}$，
$\boxed{t\simeq3.0\,\mathrm{ns}}$。
\excheck 不能用非相对论动能 $mv^2/2$ 到达 $mc^2$，因为此时粒子已明显相对论化。

\electrochargedproblem{376}{147}
一电子沿均匀电场方向做相对论运动，电场强度模为 $E$，在动能为 $T$ 时求加速度模。
\exsolution
粒子平行于受力方向时
$F=\mathrm dp/\mathrm dt=\gamma^3m_e a$。
电子动能
$T=(\gamma-1)m_ec^2$，即
$\gamma=1+T/(m_ec^2)$。
电力模 $F=eE$，故
\[
\boxed{a=\frac{eE}{m_e\gamma^3}
=\frac{eE}{m_e[1+T/(m_ec^2)]^3}}.
\]
\excheck 当 $T\to0$，恢复牛顿加速度 $eE/m_e$；高能时加速度按 $\gamma^{-3}$ 显著减小。

\electrochargedproblem{377}{148}
相对论质子以初始速度 $\bm v_0$ 在 $t=0$ 射入均匀横向电场 $\bm E\perp\bm v_0$。求：(a) 速度偏转角 $\theta(t)$；(b) 沿初始方向的速度分量 $v_x(t)$。
\exsolution
设初始方向为 x 轴、电场为 y 轴。
电场力只沿 y，因此
$p_x=p_{x0}=\gamma_0m_pv_0$，
$p_y=eEt$，
其中 $\gamma_0=(1-v_0^2/c^2)^{-1/2}$。
速度与动量共线，
\[
\boxed{\tan\theta(t)=\frac{p_y}{p_x}
=\frac{eEt}{\gamma_0m_pv_0}
=\frac{eEt}{m_pv_0}
\sqrt{1-\frac{v_0^2}{c^2}}}.
\]
粒子总能量对应
$\gamma m_p=\sqrt{m_p^2+(p_x^2+p_y^2)/c^2}$，所以
\[
\boxed{v_x(t)=
\frac{v_0}{\sqrt{1+
(1-v_0^2/c^2)(eEt/(m_pc))^2}}}.
\]
\excheck $p_x$ 守恒但 $v_x$ 不恒定，因为横向加速使 $\gamma$ 增大。

\electrochargedproblem{378}{148}
质子经 $V'=500\,\mathrm{kV}$ 加速后穿过宽度 $d=10\,\mathrm{cm}$ 的均匀横向磁场，$B=0.51\,\mathrm T$，图3.99。求离开磁场时相对初始方向的偏转角 $\alpha$。
\exsolution
非相对论动能为
$\frac12m_pv^2=eV'$，
故动量
$p=\sqrt{2m_peV'}$。
磁场中的圆轨道半径
$r=p/(eB)=\sqrt{2m_pV'/e}/B$。
从图中圆轨迹几何，跨越磁场宽度 $d$ 时
$\sin\alpha=d/r$，因此
\[
\boxed{\alpha=
\arcsin\left(dB\sqrt{\frac{e}{2m_pV'}}\right)
\simeq30^\circ}.
\]
\excheck 先检查 $d/r<1$ 才有穿过此宽度的圆轨道；500 kV 对质子的相对论修正很小。

\electrochargedproblem{379}{148}
带电粒子在均匀磁场 $B=10.0\,\mathrm{mT}$ 中沿半径 $r=100\,\mathrm{mm}$ 的圆周运动。求速度与周期，当粒子是：(a) 非相对论质子；(b) 相对论电子。
\exsolution
磁场力只改变方向，不改变动能，圆周所需动量
$p=|q|Br$。
(a) 非相对论质子
$v_p=eBr/m_p$，
$T_p=2\pi m_p/(eB)$，
\[
\boxed{v_p\simeq9.6\times10^4\,\mathrm{m/s}
\approx100\,\mathrm{km/s}},\quad
\boxed{T_p\simeq6.6\,\mathrm{\mu s}}.
\]
(b) 电子满足
$\gamma m_ev=eBr$，
即
\[
\boxed{v_e=\frac c{\sqrt{1+(m_ec/(eBr))^2}}
\simeq0.51c}.
\]
周期按几何弧长除速度：
\[
\boxed{T_e=\frac{2\pi r}{v_e}\simeq4.1\,\mathrm{ns}}.
\]
\excheck 对电子不能使用 $v=eBr/m_e$，该非相对论公式在给定数据下将不自洽。

\electrochargedproblem{380}{148}
静止质量为 $m_0$、电荷为 $q$ 的相对论粒子，在均匀磁场 $B$ 中沿半径 $r$ 的圆形轨道运动。求：(a) 动量模；(b) 动能；(c) 向心加速度模。
\exsolution
洛伦兹力等于相对论动量方向变化率：
$|q|vB=pv/r$，故
\[
\boxed{p=|q|Br}.
\]
总能量
$E_{\rm tot}=\sqrt{m_0^2c^4+p^2c^2}$，
动能
\[
\boxed{T=
m_0c^2\left[\sqrt{1+
\left(\frac{|q|Br}{m_0c}\right)^2}-1\right]}.
\]
速度
$v=pc^2/E_{\rm tot}$，
向心加速度模为
\[
\boxed{a=\frac{v^2}{r}
=\frac{c^2}{r}
\frac{(|q|Br)^2}{m_0^2c^2+(|q|Br)^2}}.
\]
\excheck $p\ll m_0c$ 时恢复 $a\simeq q^2B^2r/m_0^2$，高能时 $a<c^2/r$。

\electrochargedproblem{381}{148}
在均匀磁场中，电子与质子的圆周运动周期相比非相对论周期都增加 $\eta=1.0\%$，求二者所对应动能。
\exsolution
相对论回旋周期
$T_{\rm cyc}=2\pi\gamma m_0/(|q|B)$，
非相对论周期
$T_0=2\pi m_0/(|q|B)$，
故 $T_{\rm cyc}/T_0=\gamma=1+\eta$。
动能为
$K=(\gamma-1)m_0c^2=\eta m_0c^2$，
所以
\[
\boxed{K_e=0.01(511\,\mathrm{keV})
\simeq5.1\,\mathrm{keV}},
\]
\[
\boxed{K_p=0.01(938\,\mathrm{MeV})
\simeq9.4\,\mathrm{MeV}}.
\]
\excheck 周期增加的百分比等于 $\gamma-1$，不是速度相对光速的百分比。

\electrochargedproblem{382}{148}
电子经电压 $V=1.0\,\mathrm{kV}$ 加速，在磁场 $B=29\,\mathrm{mT}$ 中以与场线夹角 $\alpha=30^\circ$ 运动，求螺旋轨道的螺距。
\exsolution
非相对论总速度
$v=\sqrt{2eV/m_e}$，
平行磁场分量
$v_\parallel=v\cos\alpha$，
回旋周期
$T=2\pi m_e/(eB)$，
所以螺距
\[
\boxed{\ell=v_\parallel T
=\frac{2\pi\cos\alpha}{B}
\sqrt{\frac{2m_eV}{e}}
\simeq2.0\,\mathrm{cm}}.
\]
\excheck 螺距由\emph{平行}速度决定，圆周半径由\emph{垂直}速度决定，不能交换正余弦。

\electrochargedproblem{383}{148}
一束稍有发散的非相对论带电粒子经电压 $V$ 加速后从长直螺线管轴上 A 点出发。磁感应强度分别为 $B_1,B_2$ 时都在距离 $\ell$ 处出现连续两次聚焦。求比荷 $|q|/m$。
\exsolution
粒子速度约
$v=\sqrt{2|q|V/m}$，
小角度近似下轴向速度为 $v$。
每转一个回旋周期沿轴前进的螺距
\[
h(B)=v\frac{2\pi m}{|q|B}.
\]
连续聚焦值对应固定 $\ell$ 内完成整数次
$n$ 与 $n+1$ 周：
$B_n=2\pi n mv/(|q|\ell)$，
因此
\[
|B_2-B_1|=\frac{2\pi mv}{|q|\ell}.
\]
平方并代速度关系：
$(\Delta B)^2=
(4\pi^2m^2/(q^2\ell^2))(2|q|V/m)
=8\pi^2mV/(|q|\ell^2)$。
故
\[
\boxed{\frac{|q|}{m}=
\frac{8\pi^2V}{\ell^2(B_2-B_1)^2}}.
\]
\excheck 只需相邻两次聚焦的磁场差，不需要知道聚焦的具体整数 $n$。

\electrochargedproblem{384}{148}
非相对论电子从螺线管轴上 A 点以速度 $v$ 运动，速度与轴夹角 $\alpha$，均匀轴向磁场为 $B$。在距 A 轴向距离 $\ell$ 的垂直屏幕上落点距轴多远？
\exsolution
电子横向速度 $v_\perp=v\sin\alpha$，
回旋半径
$\rho=m_ev\sin\alpha/(eB)$。
轴向速度 $v_\parallel=v\cos\alpha$，
抵达屏幕时间
$t=\ell/(v\cos\alpha)$。
回旋相位
$\phi=(eB/m_e)t=eB\ell/(m_ev\cos\alpha)$。
横向圆轨道从轴上出发，初始点与到达点间弦长为
$2\rho|\sin(\phi/2)|$，故
\[
\boxed{r_{\rm hit}=
\frac{2m_ev\sin\alpha}{eB}
\left|\sin\frac{eB\ell}{2m_ev\cos\alpha}\right|}.
\]
\excheck $r_{\rm hit}$ 不超过 $2\rho$，且当相位恰为完整整数圈时落点回到轴上。

\electrochargedproblem{385}{148--149}
半径为 $a$ 的无限长直导线通电流 $I$。电子从导线外表面以初速 $v_0$ 沿径向垂直表面射出，受电流的自磁场作用转向。求电子到达的最大径向距离 $r_{\max}$。
\exsolution
导线外环向磁场
$B_\varphi(r)=\mu_0I/(2\pi r)$。
电子仅受磁力，速率恒为 $v_0$。
以 z 为导线轴线，选择磁矢势
$A_z(r)=-(\mu_0I/2\pi)\ln(r/a)$。
系统对 z 平移不变，正则动量
$p_z+qA_z$ 守恒；电子 $q=-e$，初始轴向速度为零，
故在半径 $r$ 处
\[
|mv_z|=\frac{e\mu_0I}{2\pi}\ln\frac ra.
\]
径向速度在最远点消失，且总速率为 $v_0$，
因此 $|v_z|=v_0$：
$\ln(r_{\max}/a)=2\pi m_ev_0/(e\mu_0I)$。
于是
\[
\boxed{r_{\max}=a\exp
\left(\frac{2\pi m_ev_0}{e\mu_0I}\right)}.
\]
\excheck 指数必须无量纲；导线无限长近似和无电场条件是结论前提。

\electrochargedproblem{386}{149}
非相对论带电粒子经过同轴圆柱电容器的径向电场，在电容器中沿一段圆弧运动，离开后在横向均匀磁场 $B$ 中沿半径 $r$ 的半圆运动。电容器内外电极半径为 $a<b$，加电压 $V$。求粒子速度与比荷 $q/m$。
\exsolution
圆柱电容器内径向场强模
$E(s)=V/[s\ln(b/a)]$。
粒子在电场内保持半径 $s$ 的圆弧运动，
向心平衡
$mv^2/s=|q|E(s)$，所以
\[
v^2=\frac{|q|V}{m\ln(b/a)}.
\]
在磁场内圆轨道满足
$mv^2/r=|q|vB$，即
$|q|/m=v/(Br)$。
代入上一式消去比荷，得
\[
\boxed{v=\frac{V}{Br\ln(b/a)}},
\qquad
\boxed{\frac{|q|}{m}=
\frac{V}{B^2r^2\ln(b/a)}}.
\]
\excheck 这里 $r$ 特指磁场中的半圆半径，而电容器的径向坐标 $s$ 消去；对实际带电符号应另看弯曲方向。

\electrochargedproblem{387}{149}
图3.101均匀磁场 $\bm B$ 与电场 $\bm E$ 均沿 y 轴。粒子比荷为 $q/m$，由原点以初速度 $v_0$ 沿 x 轴出发。求：(a) 第 $n$ 次回到 y 轴时纵坐标 $y_n$；(b) 当时速度与 y 轴夹角 $\alpha$。
\exsolution
磁场只影响 xz 平面速度的旋转；在 y 方向电场给
$\ddot y=qE/m$，初始 $v_y=0$，故
$y(t)=qEt^2/(2m)$。
横向回旋角频率
$\omega_c=|q|B/m$，
每完成一整圈回到 y 轴的时刻
$t_n=2\pi n/\omega_c=2\pi nm/(|q|B)$。
因此
\[
\boxed{y_n=\frac{2\pi^2n^2mE}{qB^2}}
\]
（有向 $y_n$ 的符号随 $qE$ 决定）。
第 $n$ 次回归时横向速度模仍为 $v_0$，
纵向速度
$|v_y|=|qE|t_n/m=2\pi nE/B$，
所以
\[
\boxed{\tan\alpha=\frac{v_0B}{2\pi nE}}.
\]
\excheck 电场与磁场平行，横向回旋和纵向匀加速运动可以解耦。

\electrochargedproblem{388}{149}
如图3.101一束相同比荷 $q/m$、速度不同的离子从 O 沿 x 轴射入，电场 $\bm E$ 与磁场 $\bm B$ 都沿 y 轴。垂直 x 轴的屏幕位于 $x=\ell$。求离子在屏幕上的落点轨迹方程，并说明当 $|z|\ll\ell$ 时为抛物线。
\exsolution
设正离子且 $E,B>0$，磁回旋角频率
$\omega=qB/m$。
离子在 xz 平面的轨迹为
\[
x(t)=\frac{v_0}{\omega}\sin\omega t,
\qquad
z(t)=\frac{v_0}{\omega}(1-\cos\omega t),
\]
而 y 方向有
$y(t)=qEt^2/(2m)$。
在屏幕 $x=\ell$ 上，
$z/\ell=(1-\cos\omega t)/\sin\omega t
=\tan(\omega t/2)$。
令 $\theta=\omega t/2$，则
$\theta=\arctan(z/\ell)$ 且
\[
y=\frac{qE}{2m}\left(\frac{2\theta}{\omega}\right)^2
=\frac{2mE}{qB^2}\theta^2.
\]
故屏幕痕迹的参数消去形式为
\[
\boxed{z=\ell\tan\left[
B\sqrt{\frac{qy}{2mE}}\right]}.
\]
若 $|z|\ll\ell$，$\theta\simeq z/\ell$，
则
\[
\boxed{y\simeq\frac{2mE}{qB^2\ell^2}z^2},
\]
为抛物线。
\excheck 轨迹式针对屏幕上的首次交点以及所选磁旋向；若 $q<0$，轴向偏转及相位符号需要对应修正。
